\documentclass{article}

%Chargement des paquets
\usepackage{amsmath}
\usepackage{amsfonts}
\usepackage{amssymb}
\usepackage{enumerate}
\usepackage{mathtools}
\usepackage{bbm}
\usepackage{xparse, etoolbox}
\usepackage{enumerate}
\usepackage{enumitem}
\usepackage{mathabx}
\usepackage{babel}[french]

%environment
\usepackage{amsthm}
\usepackage{keytheorems}
\usepackage{hyperref} 

%Interface théorème
\newkeytheorem{lem}[
	name = Lemme
]

\newkeytheorem{prop}[
	name = Proposition
]

\newkeytheorem{thm}[
	name = Théorème
]

\newkeytheorem{coro}[
	name = Corollaire
]

\renewcommand*{\proofname}{Démonstration}

\theoremstyle{definition}
\newtheorem{defn}{Définition}

\theoremstyle{definition}
\newtheorem*{nota}{Notation}

\theoremstyle{definition}
\newtheorem*{limi}{Liminaire}

\theoremstyle{remark}
\newtheorem{rmq}{Remarque}

\theoremstyle{plain}
\newtheorem*{fait}{Fait}


%Ensembles de nombres
\newcommand{\N} {\mathbb{N}}
\newcommand{\Ne}{\N^\ast}
\newcommand{\Z} {\mathbb{Z}}
\newcommand{\Q} {\mathbb{Q}}
\newcommand{\R} {\mathbb{R}}
\newcommand{\Rb}{\overline{\mathbb{R}}}
\newcommand{\Rp}{\R_+}
\newcommand{\Rm}{\R_-}
\newcommand{\K} {\mathbb{K}}
\newcommand{\Ket} {\mathbb{K^{\ast}}}
\newcommand{\Cx}{\mathbb{C}}

%Ensembles, tribus
\newcommand{\A}{\mathbb{A}}
\renewcommand{\P}{\mathcal{P}}
\newcommand{\T}{\mathcal{T}}
\newcommand{\F}{\mathcal{F}}
\newcommand{\C} {\mathcal{C}}
\newcommand{\ouv}{\mathcal{O}}
\newcommand{\B}{\mathcal{B}}
\newcommand{\M}{\mathcal{M}}
\newcommand{\etag}{\mathcal{E}}
\newcommand{\etagOp}{\etag^{+}(\Omega)}
\newcommand{\etagO}{\etag(\Omega)}
%%Lp avant quotientage
\newcommand{\Lic}{\mathcal{L}^1}
\newcommand{\Lpc}{\mathcal{L}^p}
\newcommand{\Linfc}{\mathcal{L}^{\infty}}
\newcommand{\Lifull}{\Lic(\Omega, \B, m)}
\newcommand{\Lpfull}{\Lpc(\Omega, \B, m)}
\newcommand{\Linffull}{\Linfc(\Omega, \B, m)}
%%Lp après quotientage
\newcommand{\Li}{L^1}
\newcommand{\Ld}{L^2}
\newcommand{\Lp}{L^p}
\newcommand{\Linf}{L^{\infty}}
\newcommand{\Listd}{\Li(\Omega, m)} 
\newcommand{\Ldstd}{\Ld(\Omega, m)} 
\newcommand{\Lpstd}{\Lp(\Omega, m)} 

%Quelques opérateurs
\DeclareMathOperator{\codim}{codim}
\DeclareMathOperator{\rg}{rg}
\DeclareMathOperator{\tr}{tr}
\DeclareMathOperator{\vect}{Vect}
\DeclareMathOperator{\ima}{Im}
\DeclareMathOperator{\id}{Id}
\DeclareMathOperator{\sign}{sign}
\DeclareMathOperator{\ext}{ext}
\newcommand{\ind}{\mathbbm{1}}
\newcommand*{\spr}[2]{\left(\,#1\,\middle|\,#2\,\right)}
\newcommand{\interior}[1]{%
  {\kern0pt#1}^{\mathrm{o}}%
}
\DeclareMathOperator{\card}{card}
\DeclareMathOperator{\ppcm}{ppcm}

%Commandes de pure flemme
\newcommand{\veps}{\varepsilon}
\newcommand{\intab}{\int_{a}^{b}}
\newcommand{\intR}{\int_{-\infty}^{\infty}}
\newcommand{\intinf}{\int_{- \infty}^{+ \infty}}
\newcommand{\equi}{\Leftrightarrow}
\newcommand{\Kev}{$\K$-espace vectoriel }
\newcommand{\Kevs}{$\K$-espaces vectoriels }
\newcommand{\Rev}{$\R$-espace vectoriel }
\newcommand{\Revs}{$\R$-espaces vectoriels }
\newcommand{\Cev}{$\Cx$-espace vectoriel }
\newcommand{\Cevs}{$\Cx$-espaces vectoriels }
\newcommand{\evn}{(E, \norm{.})}
\newcommand{\interf}[2]{[#1,#2]}
\newcommand{\limb}{\lim\limits_}
\newcommand{\lebn}{\lambda_n}
\newcommand{\pp}{\text{~presque partout}}
\newcommand{\derivp}[2]{\frac{\partial #1}{\partial #2}}
\newcommand{\famiu}{(u_i)_{i \in I}}
\newcommand{\sev}{\text{~s.e.v.}}
\newcommand{\Mnp}{M_{n,p}(\K)}
\newcommand{\Mn}{M_{n}(\K)}
\newcommand{\tq}{\text{~t.q.~}}
\newcommand{\mq}{~montrer~que~}
\newcommand{\et}{\quad \text{et} \quad}
\newcommand{\ou}{\text{~ou~}}
\newcommand{\Kx}{\K[X]}


%Commandes de gars chelou
\renewcommand{\subset}{\subseteq}

%Commande restriction, ty duckduckgo
\def\restr#1#2{\mathchoice
              {\setbox1\hbox{${\displaystyle #1}_{\scriptstyle #2}$}
              \restraux{#1}{#2}}
              {\setbox1\hbox{${\textstyle #1}_{\scriptstyle #2}$}
              \restraux{#1}{#2}}
              {\setbox1\hbox{${\scriptstyle #1}_{\scriptscriptstyle #2}$}
              \restraux{#1}{#2}}
              {\setbox1\hbox{${\scriptscriptstyle #1}_{\scriptscriptstyle #2}$}
              \restraux{#1}{#2}}}
\def\restraux#1#2{{#1\,\smash{\vrule height .8\ht1 depth .85\dp1}}_{\,#2}} 

%Quelques normes
\DeclarePairedDelimiter\abs{\lvert}{\rvert}
\DeclarePairedDelimiter\labs{\bigl\lvert}{\bigr\rvert}
\DeclarePairedDelimiter\norm{\lVert}{\rVert}
\DeclarePairedDelimiterXPP\normi[1]{}\lVert\rVert{_1}{\ifblank{#1}{\:\cdot\:}{#1}}
\DeclarePairedDelimiterXPP\normd[1]{}\lVert\rVert{_2}{\ifblank{#1}{\:\cdot\:}{#1}}
\DeclarePairedDelimiterXPP\normp[1]{}\lVert\rVert{_p}{\ifblank{#1}{\:\cdot\:}{#1}}
\DeclarePairedDelimiterXPP\normq[1]{}\lVert\rVert{_q}{\ifblank{#1}{\:\cdot\:}{#1}}
\DeclarePairedDelimiterXPP\norminf[1]{}\lVert\rVert{_\infty}{\ifblank{#1}{\:\cdot\:}{#1}}

%Bracket en angle
\DeclarePairedDelimiter\angl{\langle}{\rangle}

%Set, parenthèses
\newcommand{\set}[1]{\{ #1 \}}
\newcommand{\bigset}[1]{\bigl\{ #1 \bigr\}}
\newcommand{\Bigset}[1]{\Bigl\{ #1 \Bigr\}}
\newcommand{\bigpar}[1]{\bigl( #1 \bigr)}
\newcommand{\Bigpar}[1]{\Bigl( #1 \Bigr)}

%Opitmisation des opérateurs de comparaison
\DeclarePairedDelimiter\tio{o (}{)}
\DeclarePairedDelimiter\gro{O (}{)}


\begin{document}

\begin{thm}
Toute suite numérique croissante et majorée (resp. décroissante et minorée) est convergente.
\end{thm}

\begin{proof}
Soit $(u_n)_S$ une telle suite. L'ensemble $\set{u_n ; n \in \N}$ est non vide et majoré, il admet donc une borne supérieure.
Notons $l$ cette borne supérieure. Soit $\veps > 0$, il existe $N \in \N$ tel que $\forall n \geq N, l - \veps \leq u_n \leq l $ par définition
d'une borne supérieure, d'où le résultat.
\end{proof}

\begin{coro}
Toute suite numérique croissante (resp. décroissante et minorée) est convergente dans $\Rb$.
\end{coro}

\begin{thm}[Suites adjacentes]
Soient $(u_n)_\N$ et $(v_n)_\N$ deux suites numériques telles que :
\begin{enumerate}[label=(\roman*)]
\item $(u_n)_\N$ est croissante et $(v_n)_S$ est décroissante ;
\item la suite $(u_n-v_n)_\N$ converge vers 0.
\end{enumerate}
Alors les suites $(u_n)_\N$ et $(v_n)_\N$ sont convergentes et de même limite $l$ et de plus
\[ \forall n \in \N, u_n \leq l \leq v_n . \]
\end{thm}

\begin{proof}
Il est clair que $\forall n \in \N, u_n \leq v_n$, sinon, en prenant $N \in \N$ tel que $u_N > v_N$, on aurait, 
par monotonie des deux suites 
\[ \forall n \geq N, u_n \geq u_N > v_N \geq v_n , \]
ce qui est incompatible avec (ii).
Ainsi, la suite $(u_n)_\N$ est croissante et majorée par $v_0$ donc convergente, et de même pour $(v_n)_\N$.
L'égalité des limites est une conséquence directe de la linéarité de l'opérateur limite.
\end{proof}

\begin{thm}
De toute suite numérique, on peut extraire une suite monotone.
\end{thm}

\begin{proof}

\end{proof}

\begin{thm}[Bolzano-Weierstrass]
De toute suite numérique bornée on peut extraire une suite convergente.
\end{thm}

\begin{proof}
Soit $(u_n)_\N \subset \R$ une suite bornée. Notons $a, b \in \R$ les réels tels que $\forall n \in \N, a \leq u_n \leq b$.
On pose $a_0 = a, b_0 = b$ ainsi que les deux ensembles :
\[ \Bigset{p \in \N, a \leq u_p \leq \dfrac{a+b}2} \et \Bigset{p \in \N, \dfrac{a+b}2 \leq u_p \leq b} . \]
L'union de ces deux ensembles est égal à $\N$, donc un des deux ensembles est infini.
Si c'est le premier, on note $a_1 = a_0 ~;~ b_1 = \frac{a+b}{2}$ sinon, on note $a_1= \frac{a+b}{2} ~;~ b_1 = b_0$.
On itère ce processus et on construit ainsi deux suites $(a_n)$ et $(b_n)$ respectivement croissante et majorée et décroissante et minorée.
Par ailleurs, les intervalles $[a_n ; b_n]$ sont ainsi construits que leur diamètre est égal à $\frac{b-a}{2^n}$, donc les suites 
$(a_n)$ et $(b_n)$ sont adjacentes. On note $l$ la limite commune de ces deux suites.
On construit une extratrice $\varphi : \N \to \N$ de la façon suivante :
\begin{enumerate}[label=(\alph*)]
\item on pose $\varphi(0) = 0$ ;
\item les images de $0, 1, \dots, n$ étant supposées construites pour $n \in \N$, on construit $\varphi(n+1)$ comme suivant :
\[ u_{\varphi(n+1)} \in [a_{n+1} ; b_{n+1}] \et \varphi(n+1) > \varphi(n) , \]
ceci est toujours possible car les intervalles $[a_n; b_n]$ contiennent une infinité de points de la suite $(u_n)_\N$.
\end{enumerate}
Ainsi, par construction, pour tout $\veps > 0$, il existe $N \in \N$ tel que : 
\[ \forall n \in \N, n \geq N \Rightarrow \abs{u_{\varphi(n)} - l} \leq \veps , \]
d'où le résultat.
\end{proof}


\end{document}